\subsection{Detection error analysis}
\seclabel{analysis}
We applied the excellent detection analysis tool from Hoiem \etal
\cite{hoiem2012diagnosing} in order to reveal our
method's error modes, understand how fine-tuning changes them, and
to see how our error types compare with DPM.  A full summary of the
analysis tool is beyond the scope of this paper and we encourage readers to consult \cite{hoiem2012diagnosing} to
understand some finer details (such as ``normalized AP''). Since the
analysis is best absorbed in the context of the associated plots, we
present the discussion within the captions of \figref{fp-type-dist}
and \figref{obj-char}.
%\vspace{-0.5em}
\begin{figure}[h!]
\def \sz {1.2in}
\def \sp {-0.7em}
\centering
\includegraphics[height=\sz]{figures/det_analysis/rcnn_fc6/plots_fp_trendarea_animals.pdf}
\hspace{\sp}
\includegraphics[height=\sz,clip=true,trim=1.05cm 0 0 0]{figures/det_analysis/rcnn_ft_fc7/plots_fp_trendarea_animals.pdf}
\hspace{\sp}
%\includegraphics[height=\sz,clip=true,trim=1.05cm 0 0 0]{figures/det_analysis/dpm_v5/plots_fp_trendarea_animals.pdf}\\
\includegraphics[height=\sz,clip=true,trim=1.05cm 0 0 0]{figures/det_analysis/rcnn_ft_fc7_bb/plots_fp_trendarea_animals.pdf}\\
\vspace{2pt}
%\hspace{\sp}
\includegraphics[height=\sz]{figures/det_analysis/rcnn_fc6/plots_fp_trendarea_furniture.pdf}
\hspace{\sp}
\includegraphics[height=\sz,clip=true,trim=1.05cm 0 0 0]{figures/det_analysis/rcnn_ft_fc7/plots_fp_trendarea_furniture.pdf}
\hspace{\sp}
%\includegraphics[height=\sz,clip=true,trim=1.05cm 0 0 0]{figures/det_analysis/dpm_v5/plots_fp_trendarea_vehicles.pdf}\\
\includegraphics[height=\sz,clip=true,trim=1.05cm 0 0 0]{figures/det_analysis/rcnn_ft_fc7_bb/plots_fp_trendarea_furniture.pdf}\\
\caption{\textbf{Distribution of top-ranked false positive (FP)
types.} Each plot shows the evolving distribution of FP types as
more FPs are considered in order of decreasing score. Each FP is
categorized into 1 of 4 types: Loc---poor localization (a detection with an IoU overlap
with the correct class between 0.1 and 0.5, or a duplicate); Sim---confusion with a similar category;
Oth---confusion with a dissimilar object category; BG---a FP that fired
on background.  Compared with DPM (see \cite{hoiem2012diagnosing}), significantly more of our errors
result from poor localization, rather than confusion with background
or other object classes, indicating that the CNN features are much
more discriminative than HOG. Loose localization likely results from our
use of bottom-up region proposals and the positional invariance
learned from pre-training the CNN for whole-image classification.
Column three shows how our simple bounding-box regression method fixes many localization errors.}
\figlabel{fp-type-dist}
\vspace{-1em}
\end{figure}

\begin{figure*}[t!]
\centering
\def \sz {1.275in}
\includegraphics[height=\sz]{figures/det_analysis/rcnn_fc6/plots_impact_ov50.pdf}
\hspace{0em}
\includegraphics[height=\sz,clip=true,trim=0.5cm 0 0 0]{figures/det_analysis/rcnn_ft_fc7/plots_impact_ov50.pdf}
\hspace{0em}
\includegraphics[height=\sz,clip=true,trim=0.5cm 0 0 0]{figures/det_analysis/rcnn_ft_fc7_bb/plots_impact_ov50.pdf}
\hspace{0em}
\includegraphics[height=\sz,clip=true,trim=0.5cm 0 0 0]{figures/det_analysis/dpm_v5/plots_impact_ov50.pdf}
\caption{\textbf{Sensitivity to object characteristics.} Each plot shows the mean (over classes)
normalized AP (see \cite{hoiem2012diagnosing}) for the highest and lowest performing subsets within six different
object characteristics (occlusion, truncation, bounding-box area, aspect ratio, viewpoint, part visibility).
We show plots for our method (R-CNN) with and without fine-tuning (FT) and bounding-box regression (BB) as well as for DPM voc-release5.
Overall, fine-tuning does not reduce sensitivity (the difference between max and min), but does substantially improve
both the highest and lowest performing subsets for nearly all characteristics. This indicates that fine-tuning does
more than simply improve the lowest performing subsets for aspect ratio and bounding-box area, as one might
conjecture based on how we warp network inputs. Instead, fine-tuning improves robustness
for all characteristics including occlusion, truncation, viewpoint, and part visibility.}
\figlabel{obj-char}
\vspace{-1em}
\end{figure*}

\subsection{Bounding-box regression}
\seclabel{bboxreg}
Based on the error analysis, we implemented a simple method to reduce localization errors.
Inspired by the bounding-box regression employed in DPM \cite{lsvm-pami}, we train a linear regression model to predict a new detection window given the \pool{5} features for a selective search region proposal.
Full details are given in \asecref{bboxreg}.
Results in \tableref{voc2010}, \tableref{voc2007}, and \figref{fp-type-dist} show that this simple approach fixes a large number of mislocalized detections, boosting mAP by 3 to 4 points.

\subsection{Qualitative results}
Qualitative detection results on ILSVRC2013 are presented in \figref{examples1} and \figref{examples2} at the end of the paper.
Each image was sampled randomly from the \valb set and all detections from all detectors with a precision greater than 0.5 are shown.
Note that these are not curated and give a realistic impression of the detectors in action.
More qualitative results are presented in \figref{cexamples1} and \figref{cexamples2}, but these have been curated.
We selected each image because it contained interesting, surprising, or amusing results.
Here, also, all detections at precision greater than 0.5 are shown.
